\section{Тема 5. Векторы и матрицы}

\subsection{Постановка задачи (вариант~4)}

Куб состоит из $N^3$ прозрачных и непрозрачных элементарных
кубиков, заданных трёхмерным массивом. Построить проекцию куба
вдоль оси~$Z$: ячейка проекции видима (непрозрачна), если хотя бы
один из кубиков в соответствующем столбце вдоль оси~$Z$
непрозрачен.

\subsection{Математическое обоснование}

Пусть $\text{cube}[i][j][k] \in \{0,1\}$ --- признак непрозрачности
кубика с координатами $(i,j,k)$. Видимость ячейки $(i,j)$ проекции
на плоскость~$XY$ определяется логическим «ИЛИ» по всем значениям
$k$:
\[
  \text{proj}[i][j] = \bigvee_{k=0}^{N-1} \text{cube}[i][j][k].
\]
Общее число непрозрачных кубиков:
\[
  \text{total} = \sum_{i,j,k} \text{cube}[i][j][k].
\]

\subsection{Блок-схема алгоритма}

\begin{center}
\begin{tikzpicture}[
  node distance=10mm, start chain=going below,
  nodes={draw=black, top color=white, bottom color=lightgray!50,
         minimum width=5.6cm, align=center, font=\small},
  arrow/.style={->, >=Stealth, thick},
]
  \node[draw, rounded corners, on chain] {Start};
  \node[shape=trapezium, trapezium left angle=70,
        trapezium right angle=110, on chain, draw, join=by arrow]
       {$\mathrm{cube}[N][N][N]$};
  \node[shape=chamfered rectangle, chamfered rectangle xsep=10pt,
        on chain, draw, join=by arrow]
       (li) {$i := 0,\ N{-}1$};
  \begin{scope}[local bounding box=sci]
    \node[shape=chamfered rectangle, chamfered rectangle xsep=10pt,
          on chain, draw, join=by arrow]
         (lj) {$j := 0,\ N{-}1$};
    \begin{scope}[local bounding box=scj]
      \node[shape=rectangle, on chain, draw, join=by arrow]
           {$\mathrm{vis} := 0$};
      \node[shape=chamfered rectangle, chamfered rectangle xsep=10pt,
            on chain, draw, join=by arrow]
           (lk) {$k := 0,\ N{-}1$};
      \begin{scope}[local bounding box=sck]
        \node[shape=diamond, aspect=2, on chain, draw, join=by arrow,
              minimum width=5.0cm]
             (c) {$\mathrm{cube}[i][j][k]$?};
        \node[shape=rectangle, draw, right=2cm of c]
             (sv) {$\mathrm{vis}:=1$};
        \draw[arrow] (c.east) -- node[above,font=\scriptsize]{$+$} (sv.west);
        \coordinate[on chain, join=by arrow] (cke);
        \draw[arrow] (sv.south) |- (cke);
      \end{scope}
      \draw[arrow] (sck.south) -| ($(sck.west)-(0.9,0)$) |- (lk);
      \coordinate[on chain] (spk);
      \draw[arrow] (lk) -| ($(sck.east)+(0.9,0)$) |- (spk);
      \node[shape=trapezium, trapezium left angle=70,
            trapezium right angle=110, on chain, draw, join=by arrow]
           {$\mathrm{proj}[i][j] := \mathrm{vis}$};
      \coordinate[on chain, join=by arrow] (cje);
    \end{scope}
    \draw[arrow] (scj.south) -| ($(scj.west)-(1.1,0)$) |- (lj);
    \coordinate[on chain] (spj);
    \draw[arrow] (lj) -| ($(scj.east)+(1.1,0)$) |- (spj);
    \coordinate[on chain, join=by arrow] (cie);
  \end{scope}
  \draw[arrow] (sci.south) -| ($(sci.west)-(1.4,0)$) |- (li);
  \coordinate[on chain] (spi);
  \draw[arrow] (li) -| ($(sci.east)+(1.4,0)$) |- (spi);
  \node[draw, rounded corners, on chain, join=by arrow] {End};
\end{tikzpicture}
\end{center}

\subsection{Текст программы}

\lstinputlisting[style=Cstyle]{../src/t05.c}

\subsection{Тестовые данные}

Тестовый куб $3\times3\times3$, где непрозрачен только центральный
столбец $(1,1,*)$:

\textbf{Результат выполнения программы:}
\begin{lstlisting}[language={},basicstyle=\ttfamily\small,frame=single]
Введите куб 3x3x3 (1=непрозрачный, 0=прозрачный):
0 0 0  0 0 0  0 0 0
0 1 0  0 1 0  0 1 0
0 0 0  0 0 0  0 0 0
Проекция по оси Z:
0 0 0
0 1 0
0 0 0
Всего непрозрачных кубиков: 3 из 27
\end{lstlisting}
